% Appendix B. Implementation detail that would interrupt §4 but that a
% reader reimplementing the engine needs. No floats.

\section{Engine specification}\label{sec:appendix_b}

This appendix records three implementation details that Section~\ref{sec:method}
states but does not develop.

\subsection*{Reading the quantile off a finite scenario set}

With $\window = 1{,}260$ scenarios and $\cover = 0.05$, the 5\% point falls
between the 63rd and 64th worst outcome, and the convention for reading it is
not disclosed. Two are implemented. Linear interpolation between the bracketing
order statistics is the default. The discrete reading takes the $\lceil \cover
\window \rceil$th worst outright, which is the more conservative of the two and
what a bank wishing to avoid understating risk would choose.

The ceiling matters rather than the floor. Flooring agrees with the ceiling
whenever $\cover\window$ is fractional and silently picks the wrong side
whenever it is a whole number, which is precisely the case a unit test is
likely to be written around. Section~\ref{sec:assumptions} reports the
difference between the two conventions: 0.30\% of the VaR level and 0.06
percentage points of the mean diversification ratio.

Under age weighting the same two conventions carry over, with cumulative weight
in place of cumulative count. The weighted interpolation uses the plotting
position $(F_{j} - w_{j}) / (1 - w_{j})$, where $F_{j}$ is cumulative weight
through the $j$th ordered outcome. Under equal weights this reduces to
$(j-1)/(\window-1)$, which is what the unweighted percentile interpolates
against, so the two estimators coincide exactly at $\lambda = 1$. That identity
is what Section~\ref{sec:validation} pins to machine precision, and it is the
reason the age-weighting results in Section~\ref{sec:assumptions} can be read
as effects of the weighting rather than of a change of estimator.

\subsection*{Why the sleeve calibration is closed form}

Every position in the proxy book has profit and loss proportional to its
notional: a linear position returns the factor's return times notional, and a
bond position returns the revaluation ratio of
Equation~\eqref{eq:revaluation} times notional. Scaling a sleeve's notionals by
$k$ therefore scales every scenario outcome for that sleeve by $k$, which
scales the order statistics by $k$, which scales the sleeve's VaR by $k$ on
every date and hence its average over any period. Equation~\eqref{eq:homogeneity}
follows, and one division per sleeve lands on the disclosed target exactly.

The property is positive homogeneity, the third of the four axioms in
\citet{artzner1999coherent}. No iteration or search is involved, which is why
Table~\ref{tab:sizing} reports a residual at the scale of floating-point noise
rather than a convergence tolerance. Had the book contained options or any
other position whose value is non-linear in notional, the argument would fail
and the calibration would need solving numerically.

\subsection*{The look-ahead probe}

The timing convention --- VaR dated $t$ from scenarios up to $t$, scored
against the outcome on $t+1$ --- is enforced by a test rather than by
inspection, because look-ahead is the one error class that makes a model look
better rather than worse.

The probe multiplies every factor move after a cut-off date by 25 and re-runs
the engine on the tampered panel. Every column dated on or before the cut-off
must be identical to its untampered value. The realised profit-and-loss column
is excluded, because that column is supposed to change: it carries the outcome
of the following day, which is exactly what the tampering alters. A separate
test confirms that the realised result recorded against date $t$ is the book
repriced on the move of $t+1$ rather than of $t$, which is the off-by-one that
would make the backtest meaningless while leaving every number plausible.
